\chapter*{安全须知}
\addcontentsline{toc}{chapter}{安全须知}
\markboth{安全须知}{}

这本书相信一件事：科学是动手的学问。但动手之前，请先读完这一页。

\section*{本书实验的分级}

书中的【动手试一试】栏目只使用家庭或普通课堂可以安全获得的材料：食盐、糖、醋、小苏打、紫甘蓝、牛奶、食用油、酵母、苹果、土豆、彩笔、吸管之类。这些实验在通风良好、有成年人知晓的环境下可以完成。

以下内容\textbf{只许看、不许做}，书中只会描述或建议观看正规视频资料：

\begin{itemize}[itemsep=2pt]
\item 碱金属（锂、钠、钾）与水或酸的反应；
\item 浓酸、浓碱、氯气、漂白剂与其他清洁剂混合；
\item 放射性材料与高能辐射；
\item 有机溶剂的萃取、蒸馏，以及一切易燃蒸气与明火；
\item 细菌培养、抗生素试验和任何转基因操作；
\item 人体样本、来源不明的微生物和医学检测；
\item 强紫外线、高压电解、高温高压反应。
\end{itemize}

\section*{每条通用规则}

\begin{enumerate}[itemsep=2pt]
\item 做任何实验前，先读完整个栏目，确认材料、步骤和废物处理办法。
\item 告诉家里或学校里的成年人你在做什么；需要加热的环节由成年人操作或全程陪同。
\item 绝不品尝任何实验材料——哪怕它本来是可以吃的食物，一旦进入实验流程就不再回到餐桌。
\item 实验器具与餐具分开，用后洗净；废液按栏目说明处理（本书实验的废液基本都可用大量水冲入下水道）。
\item 不用“更好玩的替代品”：栏目说用食醋，就不要换成洁厕灵；说用小苏打，就不要换成不明的白色粉末。
\item 出了意外（打翻、溅到皮肤眼睛、误食），立即停止，用大量清水冲洗，并马上告诉成年人。
\end{enumerate}

科学史上最优秀的实验者，都是最懂得控制风险的人。谨慎不是胆小，而是专业。
